\scanpage{74}
permettant d'interpréter les éléments remarquables $\rho,\sigma_i,\tau_i,\widetilde\sigma_i,\deleted{\unsure{\tau_i}},\varepsilon_i,h_i$ de $\Pi_{03}^{D_3,\tau}$, comme des éléments de $\operatorname{Gl}(2,\mathbb Z)/\pm1$ -- que du coup on \unsure{arrive} à \unsure{remonter} de façon canonique en des éléments de $\operatorname{Gl}(2,\mathbb Z)$, donnant lieu à un formulaire des plus sympathiques, que je développe et utilise abondamment dans mes notes de l'an dernier (1981). Si on veut généraliser ceci dans le $\unsure{=|\overline m}$ sens que les iso. (40), au cas des groupes \deleted{\unsure{intermédiaires} $\Pi$} $\Pi_n$ \inserted{et variantes} associés aux \inserted{$\unsure{n}$-}polygones réguliers (le cas envisagé précédemment étant $n=3$) il me faudra trouver une interprétation des groupes analogues à $\operatorname{Sl}(2,\mathbb Z)/\pm1$ resp. $\operatorname{Gl}(2,\mathbb Z)/\pm1$, comme sous-groupes discrets \inserted{remarquables} du groupe $\operatorname{Gl}(2,\mathbb R)$...
\figure{\begin{tikzpicture}\draw (0,0)--(.65,-.85);\end{tikzpicture}}

\textcircled{5.} Dans l'interprétation ``cartographique'', les isomorphismes (40) s'explicitent comme deux équivalences de catégories, entre les cat. isotopiques des cartes ``triangulées'' \inserted{$X$}, (i.e. dont \deleted{toutes} les faces sont d'ordre divisant trois, i.e. sont des monogones ou des triangles), \deleted{avec celles}\inserted{et celle des} cartes \deleted{orientées} triangulées pondérées \inserted{$X'$}, \uline{munies} d'un
\margin{73}
